% ============================================================================
%  PARALELIZAÇÃO DA AQUISIÇÃO + LISTING DO SCRIPT (apêndice)
% ----------------------------------------------------------------------------
%  PACOTES NECESSÁRIOS no preâmbulo do main.tex:
%     \usepackage{listings}
%     \usepackage{xcolor}
%
%  Cole o bloco "CONFIG DO LISTINGS" abaixo UMA VEZ no preâmbulo (ou em um
%  arquivo de estilo). Depois use os trechos (1) e (2) onde indicado.
% ============================================================================


% ----------------------------------------------------------------------------
%  CONFIG DO LISTINGS  (colar no preâmbulo, uma única vez)
% ----------------------------------------------------------------------------
\definecolor{codegray}{gray}{0.95}
\definecolor{codecomment}{rgb}{0.25,0.50,0.35}
\definecolor{codekw}{rgb}{0.10,0.10,0.70}
\definecolor{codestr}{rgb}{0.60,0.20,0.20}

\lstdefinestyle{shell}{
  language=bash,
  backgroundcolor=\color{codegray},
  basicstyle=\ttfamily\footnotesize,
  keywordstyle=\color{codekw}\bfseries,
  commentstyle=\color{codecomment}\itshape,
  stringstyle=\color{codestr},
  numbers=left,
  numberstyle=\tiny\color{gray},
  numbersep=8pt,
  frame=single,
  rulecolor=\color{gray!50},
  breaklines=true,
  breakatwhitespace=false,
  showstringspaces=false,
  tabsize=2,
  captionpos=b,
  literate={ç}{{\c c}}1 {ã}{{\~a}}1 {á}{{\'a}}1 {é}{{\'e}}1 {í}{{\'i}}1
           {ó}{{\'o}}1 {õ}{{\~o}}1 {ú}{{\'u}}1 {â}{{\^a}}1 {ê}{{\^e}}1
}
% ----------------------------------------------------------------------------


% ============================================================================
%  (1) PARÁGRAFO PARA A SUBSEÇÃO DE PARALELIZAÇÃO (corpo do texto)
% ============================================================================

A orquestração do fluxo de trabalho --- conversão de formatos, segmentação,
simulação e cálculo de dose --- é realizada pelo Snakemake, que também registra
automaticamente as métricas de desempenho de cada etapa. A paralelização angular
da simulação no ambiente HPC, por sua vez, foi conduzida por um \textit{script}
de \textit{shell} (Apêndice~\ref{ap:lancador}), estratégia adequada à natureza
embaraçosamente paralela da aquisição, em que os processos não trocam informação
entre si. O \textit{script} percorre os ângulos iniciais de 0 a 356 em passos de
4 e lança, em segundo plano (\texttt{nohup \dots\ \&}), as 90 instâncias
independentes do simulador, cada uma processando quatro projeções ($n=4$) e
gravando sua própria saída em um arquivo de \textit{log} distinto. Dessa forma,
o sistema operacional distribui os 90 processos pelos núcleos disponíveis, e o
uso de \texttt{nohup} garante a continuidade da execução mesmo após o
encerramento da sessão de acesso remoto ao servidor.


% ============================================================================
%  (2) APÊNDICE COM O LISTING DO SCRIPT
%      (coloque após \appendix, no fim do documento)
% ============================================================================

\chapter{Script de lançamento da aquisição angular}
\label{ap:lancador}
% Se o seu modelo usar \section em vez de \chapter para apêndices, troque acima.

O \textit{script} a seguir foi utilizado no ambiente HPC para disparar, de forma
paralela, as 90 execuções independentes da simulação que compõem a aquisição
completa de $360^\circ$. Cada execução recebe como argumento o ângulo inicial
\texttt{inicio} e processa quatro projeções consecutivas.

\begin{lstlisting}[style=shell,caption={Lançador em \textit{shell} das 90
instâncias paralelas da simulação (passo angular de $4^\circ$, $n=4$ projeções
por instância).},label={lst:lancador}]
#!/bin/bash
source ~/opengate_env/bin/activate
for ((i=0; i<360; i+=4)); do
  nohup ~/opengate_env/bin/python \
    ~/Simulations_opengate/IC_Wilson/minDetec_DT_WilsonDose.py "$i" \
    > "salida${i}.log" 2>&1 &
done
\end{lstlisting}

\noindent
O laço gera os valores \texttt{0, 4, 8, \dots, 356} (90 iterações); o operador
\texttt{\&} envia cada processo para segundo plano sem aguardar seu término,
de modo que todas as instâncias são iniciadas simultaneamente. A redireção
\texttt{> "salida\$\{i\}.log" 2>\&1} consolida a saída padrão e a saída de erro
de cada instância em um arquivo de \textit{log} próprio
(\texttt{salida0.log}, \texttt{salida4.log}, \dots), facilitando o
monitoramento e a depuração individual dos \textit{jobs}.
